\chapter{Pipeline de Compilación del Firmware: de C a \texttt{.coe}}
\label{ap:pipeline_firmware}

\section{Instalación del Entorno de Compilación Cruzada}
\label{sec:instalacion_toolchain}

Todo el pipeline descripto en este apéndice depende de un compilador
cruzado capaz de generar código para la arquitectura RISC-V RV32I sin
sistema operativo de destino (\textit{bare-metal}). Este compilador no
está disponible de forma nativa en Windows, por lo que se instaló dentro
del Subsistema de Windows para Linux (WSL), ya utilizado también para
empaquetar el núcleo Ibex con FuseSoC (Sección~\ref{ssec:fuse}). La
instalación se realizó sobre una distribución Ubuntu 24.04 LTS
(\textit{noble}) corriendo bajo WSL2, mediante los paquetes oficiales del
repositorio de Ubuntu -- sin necesidad de compilar el \textit{toolchain}
desde el código fuente ni de agregar repositorios de terceros:
\begin{lstlisting}[style=bashstyle]
sudo apt update
sudo apt install gcc-riscv64-unknown-elf binutils-riscv64-unknown-elf
\end{lstlisting}

Esto instala, entre otras, las siguientes herramientas usadas a lo largo
de este apéndice, identificadas por el prefijo \texttt{riscv64-unknown-elf-}
que indica el \textit{triplet} de destino (arquitectura \texttt{riscv64},
proveedor \texttt{unknown}, sistema operativo \texttt{elf} -- es decir,
sin sistema operativo, solo el formato de archivo ejecutable):
\begin{itemize}
    \item \texttt{riscv64-unknown-elf-gcc} (versión $13.2.0$): el
    compilador cruzado en sí, invocado con las banderas
    \texttt{-march=rv32i -mabi=ilp32} para restringir la generación de
    código al subconjunto entero base de 32 bits que implementa el
    núcleo Ibex configurado.
    \item \texttt{riscv64-unknown-elf-ld} (versión $2.42$, incluida en el
    paquete \texttt{binutils-riscv64-unknown-elf}): el enlazador, invocado
    internamente por \texttt{gcc} a partir del \textit{script} de
    enlazado \texttt{link.ld} (Sección~\ref{ssec:codigo_link_ld}).
    \item \texttt{riscv64-unknown-elf-objcopy} y
    \texttt{riscv64-unknown-elf-readelf}: utilidades de manipulación e
    inspección de archivos ELF, usadas en los pasos $4$ y $5$ del pipeline
    (Sección~\ref{sec:pipeline_c_a_coe}).
    \item \texttt{riscv64-unknown-elf-objdump}: desensamblador, usado en
    el paso $7$ del pipeline como verificación final de cordura.
\end{itemize}

Se optó deliberadamente por esta vía, paquetes precompilados del
repositorio oficial de Ubuntu, en vez de compilar el \textit{toolchain}
GNU desde el código fuente (una alternativa habitual, pero que puede
demandar más de una hora de compilación y depende de una larga cadena de
bibliotecas de desarrollo), dado que Ubuntu 24.04 ya empaqueta una versión
suficientemente reciente de GCC ($13.2.0$) con soporte completo para la
extensión base de RISC-V utilizada en este trabajo. La verificación de que
el \textit{toolchain} quedó correctamente instalado y accesible se
realizó con:

\begin{lstlisting}[style=bashstyle]
riscv64-unknown-elf-gcc --version
\end{lstlisting}

confirmando que reporta \texttt{riscv64-unknown-elf-gcc
(13.2.0-11ubuntu1+12) 13.2.0} antes de proceder con la compilación de
cualquier firmware.

\section{Proceso de Conversión: de \texttt{sanity\_check.c} a \texttt{firmware.coe}}
\label{sec:pipeline_c_a_coe}

Para que el núcleo Ibex ejecute un programa, la memoria BRAM de
instrucciones debe quedar pre-cargada con su contenido binario desde el
momento en que se programa la FPGA, en el formato de inicialización de
coeficientes (\texttt{.coe}) que espera la primitiva \texttt{blk\_mem\_gen}
de Xilinx (Sección~\ref{ssec:bram_dual_port}). El camino desde el código
fuente en C hasta ese archivo \texttt{.coe} no es directo: involucra
compilar para una arquitectura sin sistema operativo (\textit{bare-metal}),
enlazar con un mapa de memoria específico, extraer el binario plano
resultante, y por último traducirlo al formato de texto que Vivado puede
interpretar como contenido inicial de la BRAM, respetando el
desplazamiento de dirección exacto en el que arranca el código.

Este proceso se automatizó en el script
\texttt{Codigos/c2coe/build\_firmware.sh}, invocado como:
\begin{lstlisting}[style=bashstyle]
./build_firmware.sh <archivo.c> <profundidad_bram_en_palabras>
\end{lstlisting}

Ejemplo real usado en este trabajo (BRAM de 4096 palabras, ver Sección~\ref{sssec:bram_dimensions}):

\begin{lstlisting}[style=bashstyle]
./build_firmware.sh sanity_check.c 4096
\end{lstlisting}

El script requiere, en el mismo directorio, dos archivos de soporte
además del \texttt{.c} a compilar: \texttt{start.S} (el punto de entrada
en ensamblador, que inicializa el puntero de pila antes de saltar a
\texttt{main}) y \texttt{link.ld} (el \textit{linker script} que fija la
dirección de origen del programa en \texttt{0x00000080}, dejando los
primeros $128$ bytes de la memoria libres por convención). El proceso
completo consta de siete pasos, cada uno con una verificación explícita
que detiene la ejecución ante cualquier inconsistencia, en vez de
generar silenciosamente un \texttt{.coe} incorrecto:

\begin{enumerate}
    \item \textbf{Verificación de archivos de entrada:} confirma que el
    \texttt{.c}, \texttt{start.S} y \texttt{link.ld} existen en el
    directorio de trabajo antes de intentar nada.

    \item \textbf{Normalización de finales de línea de \texttt{link.ld}:}
    convierte automáticamente terminaciones de línea estilo Windows
    (\texttt{CRLF}) a estilo Unix (\texttt{LF}) y elimina un eventual
    \textit{Byte Order Mark} (BOM) UTF-8 al inicio del archivo. Este paso
    es necesario porque el \textit{linker} de GNU (\texttt{ld}) puede
    fallar o interpretar mal un \textit{script} con finales de línea
    mixtos, un problema recurrente al editar archivos en un entorno
    Windows y compilarlos luego bajo WSL.

    \item \textbf{Compilación:} invoca al compilador cruzado con
    banderas fijas y conocidas para no depender de configuración
    implícita:

\begin{lstlisting}[style=bashstyle]
riscv64-unknown-elf-gcc -march=rv32i -mabi=ilp32 -nostdlib \
    -nostartfiles -O2 -T link.ld -o firmware.elf start.S archivo.c
\end{lstlisting}

    (\texttt{-march=rv32i}/\texttt{-mabi=ilp32} restringen la generación
    de código al subconjunto entero base de RV32, sin extensiones que el
    núcleo Ibex configurado no implemente; \texttt{-nostdlib}/
    \texttt{-nostartfiles} evitan enlazar contra una biblioteca estándar o
    un \textit{runtime} de arranque inexistentes en este entorno
    \textit{bare-metal}).

    \item \textbf{Verificación de la dirección real de \texttt{.text} en
    el ELF:} en vez de asumir que el código arranca en
    \texttt{0x00000080} (el valor fijado en \texttt{link.ld}), el script
    lee esa dirección directamente del archivo \texttt{.elf} recién
    generado (\texttt{readelf -S}), y aborta con un error explícito si
    tanto el punto de entrada como la dirección de \texttt{.text}
    resultan ser \texttt{0x0} -- síntoma casi seguro de que el
    \textit{linker script} no se aplicó correctamente. Este paso hace que
    el script sea robusto ante cualquier cambio futuro del mapa de
    memoria, sin necesidad de actualizar valores hardcodeados en el
    propio script.

    \item \textbf{Extracción del binario plano y generación del
    \texttt{.coe}:} \texttt{objcopy -O binary} extrae el contenido crudo
    del segmento de código y datos del \texttt{.elf}, descartando toda la
    metadata del formato ELF. Como \texttt{objcopy} omite cualquier hueco
    antes de la primera dirección con contenido real, un script auxiliar
    en Python reconstruye ese hueco inicial (rellenando con ceros desde
    \texttt{0x0} hasta la dirección real de \texttt{.text} obtenida en el
    paso anterior), agrega el relleno final hasta completar exactamente
    la profundidad de palabras indicada como segundo argumento, y escribe
    el resultado en el formato de texto que exige Vivado
    (\texttt{memory\_initialization\_radix=16;} seguido de la lista de
    palabras de 32 bits en hexadecimal, en orden \textit{little-endian}).

    \item \textbf{Verificación de tamaño contra la profundidad de la
    BRAM:} calcula la cantidad de palabras que realmente necesita el
    programa (binario más el desplazamiento inicial) y aborta si excede
    la profundidad de la BRAM indicada, en vez de generar un
    \texttt{.coe} truncado.

    \item \textbf{Desensamblado de verificación:} como último chequeo, desensambla las primeras instrucciones del \texttt{.elf}
    (\texttt{objdump -d}) para inspección visual manual, permitiendo
    confirmar que la primera instrucción real del programa aparece en la
    dirección esperada (\texttt{0x80}) y no en \texttt{0x0}.
\end{enumerate}

El archivo \texttt{firmware.coe} resultante es el mismo que referencia la
propiedad \texttt{Coe\_File} del IP \texttt{bram\_instrucciones\_v2}
(Sección~\ref{sssec:bram_dimensions}), y queda embebido directamente en el
\textit{bitstream} al sintetizar el proyecto, no requiere que el ARM
transfiera el firmware al Ibex en tiempo de ejecución. 
%(AÑADIR QUE SE IMPLEMENTARA esto lo añadire en trabajo a futuro, para hacer el ibex mas versatil) 

\section{Código Fuente de los Firmwares de Prueba}
\label{sec:codigo_firmwares}

Durante el desarrollo se utilizaron dos firmwares con propósitos
complementarios: \texttt{sanity\_check.c}, un programa mínimo de dos
escrituras fijas usado para aislar problemas de arranque de cualquier
otra lógica (Sección~\ref{ssec:efecto_sonda_nulo}), y \texttt{reporter.c},
el firmware con un bucle de $N$ iteraciones usado para ejercitar el
mecanismo de trazabilidad de forma sostenida en el tiempo
(Capítulo~\ref{ch:resultados}).

\subsection{\texttt{sanity\_check.c}}
\label{ssec:codigo_sanity\_check}

\begin{lstlisting}[language=C, basicstyle=\ttfamily\footnotesize, caption={Código fuente completo de \texttt{sanity\_check.c}.}, label={code:sanity_check.c}]
#include <stdint.h>

#define REPORTER_BASE 0x80000000
#define REG_PC (*((volatile uint32_t *) (REPORTER_BASE + 0x08)))
#define REG_D  (*((volatile uint32_t *) (REPORTER_BASE + 0x10)))

int main() {
    
    REG_D  = 0xDEADBEEF;
    REG_PC = 0xCAFEF00D;

    while (1);
    return 0;
}
\end{lstlisting}

\subsection{\texttt{reporter.c}}
\label{ssec:codigo_reporter}

\begin{lstlisting}[language=C, basicstyle=\ttfamily\footnotesize, caption={Código fuente completo de \texttt{reporter.c}, con $N=1{,}500{,}000$ (ver Sección~\ref{ssec:resultados_hardware}).}, label={code:reporter.c}]
#include <stdint.h>

// El Modulo Reporter es visto por el Ibex en la direccion 0x80000000
#define REPORTER_BASE 0x80000000

// --- Mapa de registros del Reporter v2 ---
// 0x00 / 0x04: timestamp de hardware (solo lectura, contador libre de 64 bits)
#define REG_TIME_LO   (*((volatile uint32_t *) (REPORTER_BASE + 0x00)))
#define REG_TIME_HI   (*((volatile uint32_t *) (REPORTER_BASE + 0x04)))
// 0x08 / 0x0C / 0x10: registros de evento (escritura desde el Ibex)
#define REG_PC        (*((volatile uint32_t *) (REPORTER_BASE + 0x08)))
#define REG_I         (*((volatile uint32_t *) (REPORTER_BASE + 0x0C)))
#define REG_D         (*((volatile uint32_t *) (REPORTER_BASE + 0x10)))

// Generador pseudoaleatorio ligero (LCG) para reemplazar la libreria estandar
uint32_t lcg_rand(uint32_t *state) {
    *state = *state * 1664525 + 1013904223;
    return *state;
}

int main() {
    uint32_t d = 0;
    uint32_t current_pc = 0;
    uint32_t rand_state = 12345;   // Semilla inicial
    uint32_t N = 1500000;          // ~600ms a 50MHz, calzado con la ventana de polling de recolector.py

    for (uint32_t i = 0; i < N; i++) {
        // 1. Extraer el Program Counter actual (Magia RISC-V)
        __asm__ volatile ("auipc %0, 0" : "=r" (current_pc));

        // 2. Generar numero aleatorio
        d = lcg_rand(&rand_state);

        // 3. Escribir los datos de evento; el ARM lee el timestamp de
        //    hardware (REG_TIME_LO/HI) de forma independiente via AXI
        //    cada vez que muestrea, asi que no hace falta escribirlo
        //    desde aca: es un contador libre que corre solo.
        REG_PC = current_pc;
        REG_I  = i;
        REG_D  = d;
    }

    // Bucle infinito seguro al terminar, para poder observar el estado final
    while (1);

    return 0;
}
\end{lstlisting}

\section{Componentes de Bajo Nivel: Punto de Entrada y Mapa de Memoria}
\label{sec:codigo_bajo_nivel}

Además del código en C, la compilación de cada firmware
(Sección~\ref{sec:pipeline_c_a_coe}) depende de dos archivos de soporte
que no cambian entre firmwares: el punto de entrada en ensamblador y el
\textit{script} de enlazado que define el mapa de memoria. Ambos residen
junto a \texttt{sanity\_check.c} y \texttt{reporter.c} en
\texttt{Codigos/c2coe/}.

\subsection{Punto de Entrada (\texttt{start.s})}
\label{ssec:codigo_start_s}

Como el firmware se compila con \texttt{-nostdlib -nostartfiles}
(Sección~\ref{sec:pipeline_c_a_coe}), no existe ningún \textit{runtime} de
arranque provisto por una biblioteca estándar que inicialice el puntero de
pila antes de saltar a \texttt{main}: esa responsabilidad recae
enteramente en este archivo, la primera instrucción que el Ibex ejecuta
después de reset.

\begin{lstlisting}[language={[riscv]Assembler}, basicstyle=\ttfamily\footnotesize, caption={Código fuente completo de \texttt{start.s}, el punto de entrada del firmware.}, label={code:start.s}]
.section .text.start
.global _start
_start:
    la sp, _stack_top
    call main
1:  j 1b
\end{lstlisting}

La secuencia es mínima mnemónicamente: \texttt{la sp, \_stack\_top} carga
en el registro de puntero de pila (\texttt{sp}) la dirección del símbolo
\texttt{\_stack\_top}, definido por el \textit{linker script}
(Sección~\ref{ssec:codigo_link_ld}) como el extremo superior de la RAM
disponible; \texttt{call main} salta a la función \texttt{main} del
firmware en C; y la etiqueta \texttt{1: j 1b} es una salvaguarda, un
salto incondicional a sí misma, que atrapa al núcleo en un bucle
infinito seguro en el improbable caso de que \texttt{main} retorne (los
dos firmwares de este trabajo ya terminan en su propio \texttt{while(1)},
por lo que esta salvaguarda nunca llega a ejercitarse en la práctica).


\subsection{Mapa de Memoria (\texttt{link.ld})}
\label{ssec:codigo_link_ld}

\begin{lstlisting}[language=make, basicstyle=\ttfamily\footnotesize, caption={Código fuente completo de \texttt{link.ld}, el script de enlazado que define el mapa de memoria del firmware.}, label={code:link.ld}]
MEMORY
{
  RAM (rwx) : ORIGIN = 0x00000000, LENGTH = 16K
}
SECTIONS
{
  .text 0x00000080 : { *(.text.start) *(.text*) } > RAM
  .data : { *(.data*) } > RAM
  .bss  : { *(.bss*) *(COMMON) } > RAM
  _stack_top = ORIGIN(RAM) + LENGTH(RAM);
}
\end{lstlisting}

Este \textit{script} define una única región de memoria (\texttt{RAM}) de
$16\text{K}$ (los $4096$ celdas de $32$ bits reales de la BRAM, ver
Sección~\ref{sssec:bram_dimensions}), y ubica la sección \texttt{.text}, que agrupa primero \texttt{.text.start} (el punto de entrada de
\texttt{start.s}) y luego el resto del código compilado, a partir de la
dirección \texttt{0x00000080}, dejando los primeros $128$ bytes de la
BRAM sin usar por convención. El símbolo \texttt{\_stack\_top}, referenciado
desde \texttt{start.s}, se define aquí como el extremo superior exacto de
la región \texttt{RAM} (\texttt{ORIGIN + LENGTH}), de modo que la pila
crece hacia abajo desde el final de la memoria disponible, tan lejos como
sea posible del código y los datos del programa.

\section{Script Alternativo de Conversión: \texttt{bin2coe.py}}
\label{sec:bin2coe}

El paso 5 de \texttt{build\_firmware.sh} (Sección~\ref{sec:pipeline_c_a_coe})
resuelve la conversión de binario plano a \texttt{.coe} mediante un
fragmento de Python embebido directamente en el propio \textit{script} de
\textit{shell}. Ese fragmento no es más que una copia funcionalmente
equivalente de \texttt{Codigos/c2coe/bin2coe.py}, un script independiente
y reutilizable que implementa la misma lógica de relleno (\textit{padding})
como una función invocable por separado, útil para regenerar o depurar un
\texttt{.coe} a partir de un \texttt{.bin} ya existente sin necesidad de
volver a compilar ni de contar con \texttt{start.s}/\texttt{link.ld} a
mano:

\begin{lstlisting}[language=Python, basicstyle=\ttfamily\footnotesize, caption={Código fuente completo de \texttt{bin2coe.py}.}, label={code:bin2coe.py}]
import sys

def bin2coe(bin_path, coe_path, depth_words, start_addr=0x80, word_size=4, radix=16):
    """
    Convierte un binario plano (extraído con objcopy -O binary) a .coe,
    compensando el hecho de que objcopy omite cualquier hueco antes de la
    primera dirección con contenido real (en nuestro caso, 0x80, definido
    por `. = 0x00000080;` en link.ld).

    Sin este padding, el primer byte del .bin queda alineado a la
    dirección 0x00 de la BRAM en vez de 0x80, corriendo todo el programa
    hacia atrás y haciendo que el Ibex arranque en medio de una
    instrucción equivocada.
    """
    with open(bin_path, "rb") as f:
        data = f.read()

    # 1. Reconstruir el hueco inicial (0x00 .. start_addr-1) con ceros
    padded = b'\x00' * start_addr + data

    total_bytes = depth_words * word_size
    if len(padded) > total_bytes:
        raise ValueError(
            f"El binario con padding ({len(padded)} bytes, incluyendo "
            f"{start_addr} bytes de offset) no entra en la BRAM "
            f"({total_bytes} bytes). Programa demasiado grande o "
            f"profundidad de BRAM insuficiente."
        )

    # 2. Padding final hasta completar la profundidad exacta de la BRAM
    padded += b'\x00' * (total_bytes - len(padded))

    words = []
    for i in range(0, len(padded), word_size):
        val = int.from_bytes(padded[i:i+word_size], byteorder="little")
        words.append(f"{val:0{word_size*2}X}")

    with open(coe_path, "w") as f:
        f.write(f"memory_initialization_radix={radix};\n")
        f.write("memory_initialization_vector=\n")
        f.write(",\n".join(words))
        f.write(";\n")

    print(f"OK: {coe_path} generado. {len(words)} palabras totales "
          f"({start_addr} bytes de offset inicial + {len(data)} bytes de "
          f"programa real + padding final).")

if __name__ == "__main__":
    if len(sys.argv) < 4:
        print("Uso: python3 bin2coe.py <entrada.bin> <salida.coe> <profundidad_en_palabras> [start_addr_hex]")
        print("Ejemplo: python3 bin2coe.py firmware.bin firmware.coe 2048 0x80")
        sys.exit(1)
    start = int(sys.argv[4], 16) if len(sys.argv) > 4 else 0x80
    bin2coe(sys.argv[1], sys.argv[2], depth_words=int(sys.argv[3]), start_addr=start)
\end{lstlisting}

Invocado de forma independiente (por ejemplo, a partir de un
\texttt{firmware.bin} ya generado por otra vía), su uso es:
\begin{lstlisting}[style=bashstyle]
    python3 bin2coe.py firmware.bin firmware.coe 4096 0x80
\end{lstlisting}
donde el tercer argumento es la profundidad de la BRAM en palabras de 32
bits, y el cuarto (opcional, \texttt{0x80} por defecto) es la dirección
de inicio del código -- debe coincidir siempre con la dirección fijada en
\texttt{link.ld} (Sección~\ref{ssec:codigo_link_ld}).
